\DocumentMetadata{lang=he}
\documentclass[a4paper]{article}

\usepackage[hebrew, bidi=basic, layout=counters lists footnotes captions
extras sectioning.tabular, provide={onchar=ids fonts, alph=letters.plain, Alph=letters}]{babel}
\babelprovide[onchar=fonts letters]{english}

\directlua{
  luaotfload.add_fallback("hebrewrm", {
    "David CLM:mode=harf",
    "Hadasim CLM:mode=harf",
    "Liberation Serif:mode=harf",
  })
}

\babelfont[hebrew, Ligatures=TeX]{rm}[RawFeature={fallback=hebrewrm}]{David CLM}
\babelfont[hebrew, Ligatures=TeX]{sf}{Nachlieli CLM}
\babelfont[hebrew]{tt}{Liberation Mono}

\usepackage{enumitem}
\usepackage{graphicx}
\usepackage[colorlinks]{hyperref}

\NewDocumentCommand\downloadfile{m m}{%
  \IfFileExists{#1}{}{\directlua{
      os.execute([[mkdir -p "$(dirname #1)"]])
      os.execute([[curl -Lo #1 #2]])
  }}%
}

\setlocalecaption{hebrew}{figureautoref}{איור}
\setlocalecaption{hebrew}{sectionautoref}{חלק}

\title{דו״ח פרויקט: זיהוי קול}
\author{יונתן דפנא}

\begin{document}
\maketitle

\begin{abstract}
  נראה לי שלא באמת צריך תקציר.
\end{abstract}

\section{מטרת הפרויקט (מה הפרויקט עושה?)}
\label{sec:whatisit}

המטרה היא אותה המטרה כמו באבן דרך 1: לבנות מערכת שיודעת (באמצעות
contrastive learning) להבחין אם שתי הקלטות קול קצרות הוקלטו ע״י אותו
דובר.

האפליקציה הסופית יכולה לתחזק מסד נתונים של דוברים ידועים, ולהשוות
הקלטות חדשות אליהם, אבל זה לא חובה. העיקר הוא להשוות בין שתי הקלטות
נתונות.

\section{מקור הנתונים}
\label{sec:dataset}

החלטתי להשתמש ב-\href{https://www.openslr.org/12}{LibriSpeech} כמקור
נתונים. ניתן לקרוא עליו כאן:
\url{https://www.danielpovey.com/files/2015_icassp_librispeech.pdf}.

הנתונים כבר מחולקים מראש ל-train, dev ו-test, אז אפשר פשוט להשתמש
בחלוקה הקיימת. תוך כדי פיתוח (ובשביל איטרציה מהירה יותר) אפשר לפתח על
החלוקות שנחשבות ”clean” (אודיו ללא הפרעות), כי הן קטנות יותר, ובסופו
של דבר לרוץ על כל הנתונים.

LibriSpeech מבוסס על \href{https://librivox.org/}{LibriVox}, פרויקט
להקלטת ספרים בחינם, ומחלק כל פרק למספר קטעים קצרים (כל אחד כמה מילים).
על כן, הוא נתון בהיררכיה של דובר-פרק-הקלטה (ובא עם קבצי CSV שמתאימים
כל פרק לשמו ולספר).

בפרויקט הזה אין צורך להתחשב בפרק או בספר, דרושות רק הקלטות קצרות. על
כן שגרת הורדת הנתונים בקוד שלי משטחת לחלוטין את היררכית הפרקים.

\section{מודלים}
\label{sec:models}

יש שלושה רעיונות כלליים למימוש הפרויקט, ולכל אחד מהם מתאים מודל אחד או
יותר. לכולם אותו מבנה ברמת המאקרו: חזרה על התבנית
קונבולוציה-נורמליזציה-MaxPool עם כמה שכבות Dense בסוף, שפולטת וקטור
embedding באורך מסוים שניתן להשוות עם וקטורים אחרים.

\subsection{רשת קונבולוציה חד-ממדית על אודיו}
\label{sec:raw1d}

הרעיון כאן הוא לקחת את המידע בקבצי האודיו כמו שהוא (ניתן לחשוב עליו
כצורת גל של ממש, ראו \autoref{fig:waveform}), ומאחר שהוא קיים במימד
אחד (הזמן), להעביר אותו דרך מספר שכבות קונבולוציה חד-ממדיות וכמה שכבות
FC בסוף. המודל הזה \emph{לא כולל} ביצוע DFT. ניתן לדמיין את המודל כפס
אנכי שעובר על גבי צורת הגל.

\begin{figure}
  \centering
  \downloadfile{images/waveform.png}{https://forum.audacityteam.org/uploads/default/original/3X/8/1/816771cc4529dce39976d77c4150a05947a61e83.png}
  \includegraphics[scale=0.2]{images/waveform.png}
  \caption{צורת גל}
  \label{fig:waveform}
\end{figure}

\subsection{רשת קונבולוציה חד-ממדית על ספקטרוגרמה}
\label{sec:dft1d}

כאן מבצעים סוג של DFT (למעשה יוצרים Mel Spectrogram) שממיר את נתוני
האודיו ממימד אחד, לחץ אוויר מול \emph{זמן}, לשני מימדים, אמפליטודה מול
\emph{זמן ותדר} (\autoref{fig:spectrogram}). האינטואיציה היא שההבדלים
בצורת \href{https://en.wikipedia.org/wiki/Vocal_tract}{מערכת הקול} של
בני אדם שונים משפיעים על תדרים באופן יותר ישיר מעל צורת הגל.

הספקטרוגרמה היא דו-ממדית, אבל המודל הזה עדיין מבצע קונבולוציה
חד-ממדית. הרעיון הגדול מאחורי רשתות הקונבולוציה הוא שהקונבולוציה
מסייעת לשמור על המבנה המרחבי של הפרמטרים ובכך מובילה להפקת לקחים טובה
יותר. ואמנם תדירויות קרובות (הממד השני שהוספנו) באמת נשמעות לאוזנינו
קרובות, אבל, בהתייחסות לתדרים, הגורם שהכי משפיע על
\href{https://he.wikipedia.org/wiki/גוון_(מוזיקה)}{גוון הקול} הוא
דווקא היחס בין ה\emph{אוברטונים} (או \emph{הרמוניות}), שיכולים להיות
די רחוקים אחד מן השני בתדר.

על כן עדיין נבצע קונבולוציה חד-ממדית (רצה במימד הזמן), על מנת ”לראות”
את כל תחום התדרים ביחד. יש לציין שהלקחים האמיתיים מופקים רק מניסיון
ובדיקה (ראו \autoref{sec:performance}), ולכן הפרויקט בודק מספר מודלים
שונים.

\subsection{רשת קונבולוציה דו-ממדית על ספקטרוגרמה}
\label{sec:dft2d}

הרעיון דומה למודל הקודם, אבל כאן רשת הקונבולוציה היא דו-ממדית כפי
שנהוג ב-Machine Vision. אין הרבה מה לחדש על רשתות קונבולוציה דו-ממדיות
לסיווג/אימות, זה נושא שנחקר כבר רבות. בפרויקט הזה ממומשים שלושה מודלים
ע״פ הרעיון הזה: רשת קונבולוציה חדשה לחלוטין, רשת קיימת
(\href{https://arxiv.org/abs/1704.04861}{MobileNet}) מאומנת מחדש על
ספקטרוגרמות, ורשת מאומנת מראש (MobileNet על
\href{https://doi.org/10.1167\%2F9.8.1037}{ImageNet}) עם Transfer
Learning.

\begin{figure}
  \centering
  \downloadfile{images/spectrogram.png}{https://upload.wikimedia.org/wikipedia/commons/c/c5/Spectrogram-19thC.png}
  \includegraphics[scale=0.4]{images/spectrogram.png}
  \caption{דוגמה לספקטרוגרמה}
  \label{fig:spectrogram}
\end{figure}

\section{הסבר קוד}
\label{sec:code}

את הקוד ניתן לראות בקישור:
\url{https://github.com/Dafna3316/speaker-identification} (כולל קוד
המקור של הדו״ח הזה!).
אציין בחלק הזה שהרצת הקוד דורשת את התוכנות והחבילות הבאות:
\begin{itemize}
\item פייתון מעל 3.11
\item Tensorflow (גרסה 2.17 על macOS, או כל גרסה חדשה יותר על מערכות אחרות)
\item על \foreignlanguage{english}{macOS}: \foreignlanguage{english}{\texttt{tensorflow-metal}}
\item \texttt{ipywidgets} ו-\texttt{jupyterlab}
\item מערכת מבוססת UNIX® או דמוית UNIX®, כגון macOS או GNU, הכוללת: \texttt{tar}, \texttt{grep} ו-\texttt{gunzip}
\item cURL
\item FFmpeg
\end{itemize}

המחשב שהשתמשתי בו להריץ את הקוד הוא Mac mini 2024 עם M4, שיש לו 10
ליבות CPU ו-10 ליבות GPU, עם זיכרון משותף ביניהם. במעבד יש גם רכיב
מיוחד להרצת רשתות נוירונים, ה-\foreignlanguage{english}{Neural Engine
  (NE)}, אבל הוא מיועד ל-Inference ולכן לא עוזר באימון.

\subsection{הורדת נתונים}
\label{sec:code-dataset}

החלק הראשון בקוד אחראי להוריד ולעבד את הנתונים לפורמט ש-Tensorflow
יודע לעבוד איתו. בהתחלה מוגדרים כמה משתנים שקובעים אילו נתונים מתוך
LibriSpeech יש להוריד, אחר כך מחלצים את הקבצים, ולבסוף בונים את אוביקט
ה-\texttt{Dataset}.

\subsubsection{מילה על פורמט אודיו}
\label{sec:wav}

הקבצים ב-LibriSpeech נתונים בפורמט FLAC, אותו Tensorflow לא יודעת
לקרוא. קיימת חבילה \texttt{tensorflow-io} שתומכת בקריאת הפורמט הזה
(והרבה אחרים) וגם בכל מיני טרנספורמציות על אודיו, אבל היא לא עובדת על
macOS/ARM (נכשלת עם שגיאות).

הנושא הזה גרם לי הרבה כאב ראש; בהתחלה רציתי להוריד את הנתונים
מ-\href{https://www.tensorflow.org/datasets/catalog/librispeech}{TFDS},
אבל ההורדה לקחה יותר מדי זמן. שימוש ב-\texttt{tensorflow-io} קיצר את
ההורדה אבל לא הצליח לפתוח את הקבצים, ובסופו של דבר נאלצתי להוריד את
הנתונים ידנית ולהמיר אותם לפורמט WAV.

\subsubsection{הפונקציה \texttt{ensure\_set}}
\label{sec:code-ensure_set}

הפונקציה הזו מורידה קובץ מכווץ יחיד מתוך LibriSpeech ומחלצת אותו. היא
משתמש בפקודות מערכת UNIX ולכן לא תעבוד על Windows.

בקבצים המסופקים ע״י LibriSpeech, כל תיקיית דובר מכילה תיקיות פרקים.
הפונקציה הזו דואגת גם לחלץ את כל הקבצים של אותו דובר לתיקייה אחת,
ולהמיר אותם לפורמט WAV.

\subsubsection{בניית ה-\texttt{Dataset}}
\label{sec:code-dataset-obj}

הפונקציה \texttt{ds\_for\_siamese\_network} יוצרת \texttt{Dataset}
שמורכב מזוגות חיוביים ושליליים, ומתאים להעביר דרך רשת סיאמית.

היא בוחרת רק את הדוברים שיש להם יותר משתי הקלטות, ומייצרת את כל הזוגות
האפשריים של הקלטות (ומערבבת). כמות ההקלטות מוגבלת ע״י פרמטר. לכל זוג
חיובי כזה, הפונקציה מייצרת גם זוג שלילי (זוג של הקלטה מאותו דובר
והקלטה מדובר אחר), וחותכת/מרפדת את כל ההקלטות לאורך הנתון (כרגע שתי
שניות).

ככה נבנים הנתונים עבור train, test ו-dev.

\section{הערכת ביצועים}
\label{sec:performance}

\section{ממשק למשתמש}
\label{sec:app}


\section{אפשרויות לשיפור}
\label{sec:improvements}

הצעד הבא שאני מתכנן הוא לנסות להדר את \texttt{tensorflow-io} על המחשב.
אם אצליח, זה יאפשר לי להשתמש בהרבה יותר נתונים (להמיר את הכול ל-WAV
לוקח זמן) וכך לשפר את ביצועי כל המודלים.

\end{document}
